\noindent\textbf{Elección del procedimiento.} La ley de Sutherland dada es una correlación: no se deduce de este enunciado. Primero se forma el cociente absoluto $T/T_0$, se eleva a $3/2$ y se aplica el factor $(T_0+S)/(T+S)$. Se repite el mismo procedimiento a cada temperatura.

\begin{enumerate}

\item Usar 288 K para la comparación pedida, conservando $T_0=273.15$ K.

\[\fitmath{\mu_{288}=\frac{\mu_{0}\cdot {\left(\frac{\mathrm{Tcmp}}{T_{0}}\right)}^{1.5}\cdot \left(T_{0}+S\right)}{\mathrm{Tcmp}+S}}\]

\[\fitmath{\mu_{288}=\frac{(1.716\times10^{-5})\cdot {\left(\frac{(288)}{(273.15)}\right)}^{1.5}\cdot \left((273.15)+(110.4)\right)}{(288)+(110.4)}\approx 1.788574\times10^{-5}\,\mathrm{Pa\,s}}\]

\item Aplicar la misma relación a 1000 K.

\[\fitmath{\mathrm{muT}=\frac{\mu_{0}\cdot {\left(\frac{T}{T_{0}}\right)}^{1.5}\cdot \left(T_{0}+S\right)}{T+S}}\]

\[\fitmath{\mathrm{muT}=\frac{(1.716\times10^{-5})\cdot {\left(\frac{(1000)}{(273.15)}\right)}^{1.5}\cdot \left((273.15)+(110.4)\right)}{(1000)+(110.4)}\approx 4.152006\times10^{-5}\,\mathrm{Pa\,s}}\]

\end{enumerate}
